\section{AIME II 2004}
        \begin{enumerate}
        	\item (\textit{AIME II 2004}) A chord of a circle is perpendicular to a radius at the midpoint of the radius. The ratio of the area of the larger of the two regions into which the chord divides the circle to the smaller can be expressed in the form $\frac{a\pi+b\sqrt{c}}{d\pi-e\sqrt{f}},$ where $a, b, c, d, e,$ and $f$ are positive integers, $a$ and $e$ are relatively prime, and neither $c$ nor $f$ is divisible by the square of any prime. Find the remainder when the product $abcdef$ is divided by 1000.


 
 \textbf{Answer: }592



\textbf{Solution}

Let $r$ be the \href{https://artofproblemsolving.com/wiki/index.php?title=Length}{length} of the radius of the circle.  A \href{https://artofproblemsolving.com/wiki/index.php?title=Right\_triangle}{right triangle} is formed by half of the chord, half of the radius (since the chord \href{https://artofproblemsolving.com/wiki/index.php?title=Bisect}{bisects} it), and the radius. Thus, it is a $30^\circ$ - $60^\circ$ - $90^\circ$ \href{https://artofproblemsolving.com/wiki/index.php?title=Triangle}{triangle}, and the area of two such triangles is $2 \cdot \frac{1}{2} \cdot \frac{r}{2} \cdot \frac{r\sqrt{3}}{2} = \frac{r^2\sqrt{3}}{4}$. The \href{https://artofproblemsolving.com/wiki/index.php?title=Central\_angle}{central angle} which contains the entire chord is $60 \cdot 2 = 120$ \href{https://artofproblemsolving.com/wiki/index.php?title=Degree}{degrees}, so the area of the \href{https://artofproblemsolving.com/wiki/index.php?title=Sector}{sector} is $\frac{1}{3}r^2\pi$; the rest of the area of the circle is then equal to $\frac{2}{3}r^2\pi$.

 The smaller area cut off by the chord is equal to the area of the sector minus the area of the triangle. The larger area is equal to the area of the circle not within the sector and the area of the triangle. Thus, the desired ratio is $\frac{\frac{2}{3}r^2\pi + \frac{r^2\sqrt{3}}{4}}{\frac{1}{3}r^2\pi - \frac{r^2\sqrt{3}}{4}} =  \frac{8\pi + 3\sqrt{3}}{4\pi - 3\sqrt{3}}$

 Therefore, $abcdef = 2^53^4 = 2592 \Longrightarrow \boxed{592}$.

 




	\item (\textit{AIME II 2004}) A jar has 10 red candies and 10 blue candies. Terry picks two candies at random, then Mary picks two of the remaining candies at random. Given that the probability that they get the same color combination, irrespective of order, is $m/n,$ where $m$ and $n$ are relatively prime positive integers, find $m+n.$


 
 \textbf{Answer: }441



\textbf{Solution}

The probability that Terry picks two red candies is $\frac{10 \cdot 9}{20 \cdot 19} = \frac{9}{38}$, and the probability that Mary picks two red candies after Terry chooses two red candies is $\frac{7\cdot8}{18\cdot17} = \frac{28}{153}$.  So the probability that they both pick two red candies is $\frac{9}{38} \cdot \frac{28}{153} = \frac{14}{323}$.  The same calculation works for the blue candies.  

 The probability that Terry picks two different candies is $\frac{20\cdot10}{20\cdot19} = \frac{10}{19}$, and the probability that Mary picks two different candies after Terry picks two different candies is $\frac{18\cdot 9}{18\cdot 17} = \frac{9}{17}$. Thus, the probability that they both choose two different candies is $\frac{10}{19}\cdot\frac{9}{17} = \frac{90}{323}$.  Then the total probability is

 
 \[2 \cdot \frac{14}{323} + \frac{90}{323} = \frac{118}{323}\]

 and so the answer is $118 + 323 = \boxed{441}$.

 
In the above calculations, we treated the choices as ordered; that is, Terry chose first one candy, then a second, and so on.  We could also solve the problem using unordered choices.  The probabilities calculated will all be the same, but the calculations will appear somewhat different.  For instance, the probability that Mary chooses two red candies after Terry chose two red candies will have the form $\frac{{8\choose 2}}{{18 \choose 2}}$, and the probability that Terry chooses two different candies will have the form $\frac{{10\choose 1}\cdot{10\choose 1}}{{20\choose2}}$.  It is not difficult to see that these yield the same results as our calculations above, as we would expect.

 




	\item (\textit{AIME II 2004}) A solid rectangular block is formed by gluing together $N$ congruent 1-cm cubes face to face. When the block is viewed so that three of its faces are visible, exactly 231 of the 1-cm cubes cannot be seen. Find the smallest possible value of $N.$


 
 \textbf{Answer: }384



\textbf{Solution}

The $231$ cubes which are not visible must lie below exactly one layer of cubes.  Thus, they form a rectangular solid which is one unit shorter in each dimension.  If the original block has dimensions $l \times m \times n$, we must have $(l - 1)\times(m-1) \times(n - 1) = 231$.  The \href{https://artofproblemsolving.com/wiki/index.php?title=Prime\_factorization}{prime factorization} of $231 = 3\cdot7\cdot11$, so we have a variety of possibilities; for instance, $l - 1 = 1$ and $m - 1 = 11$ and $n - 1 = 3 \cdot 7$, among others.  However, it should be fairly clear that the way to minimize $l\cdot m\cdot n$ is to make $l$ and $m$ and $n$ as close together as possible, which occurs when the smaller block is $3 \times 7 \times 11$.  Then the extra layer makes the entire block $4\times8\times12$, and $N= \boxed{384}$.

 An alternate way to visualize the problem is to count the blocks that can be seen and subtract the blocks that cannot be seen.  In the given block with dimensions $l\times m \times n$, the three faces have $lm$, $mn$, and $ln$ blocks each.  However, $l$ blocks along the first edge, $m$ blocks along the second edge, and $n$ blocks along the third edge were counted twice, so they must be subtracted.  After subtracting these three edges, 1 block has not been counted - it was added three times on each face, but subtracted three times on each side.  Thus, the total number of visible cubes is $lm+mn+ln-l-m-n+1$, and the total number of invisible cubes is $lmn-lm-mn-ln+l+m+n-1$, which can be factored into $(l-1)(m-1)(n-1)$.

 




	\item (\textit{AIME II 2004}) How many positive integers less than 10,000 have at most two different digits?


 
 \textbf{Answer: }927



\textbf{Solution}

First, let's count numbers with only a single digit.  We have nine of these for each length, and four lengths, so 36 total numbers.

 Now, let's count those with two distinct digits.  We handle the cases "0 included" and "0 not included" separately.

 There are ${9 \choose 2}$ ways to choose two digits, $A$ and $B$.  Given two digits, there are $2^n - 2$ ways to arrange them in an $n$-digit number, for a total of $(2^1 - 2) + (2^2 - 2) + (2^3 -2) + (2^4 - 2) = 22$ such numbers (or we can list them: $AB, BA, AAB, ABA, BAA, ABB, BAB, BBA, AAAB, AABA, ABAA,$ $BAAA, AABB, ABAB, BAAB, ABBA, BABA, BBAA, ABBB, BABB, BBAB, BBBA$).  Thus, we have ${9 \choose 2} \cdot 22 = 36\cdot22 = 792$ numbers of this form.

 Now, suppose 0 is one of our digits.  We have nine choices for the other digit.  For each choice, we have $2^{n - 1} - 1$ $n$-digit numbers we can form, for a total of $(2^0 - 1) + (2^1 - 1) + (2^2 - 1) + (2^3 - 1) = 11$ such numbers (or we can list them: $A0, A00, A0A, AA0, A000, AA00, A0A0, A00A, AAA0, AA0A, A0AA$).  This gives us $9\cdot 11 = 99$ numbers of this form.

 Thus, in total, we have $36 + 792 + 99 = \boxed{927}$ such numbers.

 





\textbf{Solution 2a}

We use casework on the number of digits for this problem.

 If the number has a single digit, namely the number $n \in [1,9],$ we can clearly all such $n$ work.

 If the number has two digits, or the number $n \in [10,99]$ we can clearly see all such $n$ work.

 If the number $n$ has three digits, there are a total of $900$ three digit numbers, and there are $9 \cdot 9 \cdot 8$ numbers that have all distinct digits so there are $900 - 9 \cdot 9 \cdot 8$ total three digit numbers that work.

 If the number $n$ has four digits, then we have a total of $9 +\binom{9}{2}\left(\binom{4}{1}+\binom{4}{2} +\binom{4}{3}\right) + 9 \cdot \left(\binom{3}{1}+\binom{3}{2} +\binom{3}{3}\right)$ so in total we get $\boxed{927}$ numbers that work.

 
Alternatively, $n$ has four digits can be calculated as follows: When $n$ has one unique digit, there are $9$ possibilities. When $n$ has two unique digits there are two cases. Case 1: two digits are the same with each other and different with the other pair of similar digits. In that case there are $\frac{\binom{4}{2}(9^2)}{2}$ numbers that work. The case had to be divided by $2$ because it counts the cases of distinct digits $(x,y)$ again when $(y,x)$ is picked. Case 2: three digits are the same and one is different. There are $4(9^2)$ numbers that work.

 





\textbf{Solution 2b}

With subcases on the individual digits

 1 digit: $9$ ways as above.

 2 digit: $90$ ways as above.

 3 digit: $9(1\cdot10+9\cdot2)=252$ because 1st digit all except $0$ (9 ways). If 2nd digit same as 1st (1 way), 3rd can be anything (10 ways). If 2nd different than 1st (9 including 0), 3rd can be either 2nd or 1st (2). 

 4 digit: $9(1(1\cdot10+9\cdot2)+9\cdot2\cdot2)=576$ in the same way as before. 

 $9+90+252+576 = \boxed{927}$.

 





\textbf{Solution 3}

We go with complementary counting. 3 different digits. First, we need atleast a three digit number. There are $9 \cdot 9 \cdot 8$ = 648 ways to do this with a three digit number. Note that I can simply add any one of the digits from $1-9$ to the left of any of these digits to get a 4 digit number, that has atleast 3 distinct digits (It can be either of some form like ABCA or ABCD). Thus there will be $9 \cdot 648 = 5832$ such numbers. But wait, we have missed all the cases where the third digit is $0$. There are $9 \cdot 9 \cdot 8 = 648$ such ways (There are two $9's$ there because I can repeat atleast one digit). This count came to me while typing. Originally, I just did it regularly. Take the case of all 4 distinct, then the case of two repeating where one of them is the left-most digit. Now, we are missing the cases where two digits from the last three places repeat in a 4 digit number. All the cases where the first digit repeats with any of the last 3, and the cases where all 4 are distinct have been covered. For this, we split it into two cases. First case is such that the repeating digit is not $0$. We have $\binom{3}{2} \cdot \binom{9}{1} \cdot 8 \cdot 8 = 1728$ such cases. Now, we will let $0$ be the repeating digits to get $\binom{3}{2} \cdot 9 \cdot 8 = 216$ cases.
 \[648 + 5832 + 648 + 1728 + 216 = 9072 \implies \boxed{9999 - 9072 =927}\]The question says $less$ than $10,000$, so we subtract from $9999$, not $10,000$.

 




	\item (\textit{AIME II 2004}) In order to complete a large job, 1000 workers were hired, just enough to complete the job on schedule. All the workers stayed on the job while the first quarter of the work was done, so the first quarter of the work was completed on schedule. Then 100 workers were laid off, so the second quarter of the work was completed behind schedule. Then an additional 100 workers were laid off, so the third quarter of the work was completed still further behind schedule. Given that all workers work at the same rate, what is the minimum number of additional workers, beyond the 800 workers still on the job at the end of the third quarter, that must be hired after three-quarters of the work has been completed so that the entire project can be completed on schedule or before?


 
 \textbf{Answer: }766



\textbf{Solution 1}

A train is traveling at $1000$ miles per hour and has one hour to reach its destination $1000$ miles away. After $15$ minutes and $250$ miles it slows to $900$ mph, and thus takes $\frac{250}{900}(60)=\frac{50}{3}$ minutes to travel the next $250$ miles. Then it slows to $800$ mph, so the next quarter takes $\frac{250}{800}(60)=\frac{150}{8}$. The train then has $60-15-\frac{50}{3}-\frac{150}{8}=230/24$ minutes left to travel 250 miles, and doing the arithmetic shows that during this last stretch it must travel more than $1565$ mph to make it on time. Therefore the company needs to add $1566-800 = \boxed{766}$ more workers.

 Solution by rocketscience

 





\textbf{Solution 2}

Let each worker's speed be $w$, the entire time be $t$, and the total work be $1$.

 
From the initial problem statement, we have $1000w\cdot\frac{1}{4}t=\frac{1}{4}$.

 
We need to find the time the next quarter takes to complete the same amount of work, which is $\frac{1000}{900}\cdot\frac{1}{4}\cdot t=\frac{5}{18}t$.

 
Similarly, we find that the time the third quarter takes is $\frac{1000}{800}\cdot\frac{1}{4}\cdot t=\frac{5}{16}t$.

 
Finally, the time the last quarter takes is $t\left[1-\left(\frac{1}{4}+\frac{5}{18}+\frac{5}{16}\right)\right]=\frac{23}{144}t$.

 
We then let the number of workers needed be $x$, so we have the equation $\left(800+x\right)w\cdot\frac{23}{144}t=\frac{1}{4}$. Dividing by the first equation, we have

 
 \[\frac{800+x}{1000}\cdot\frac{\frac{23}{144}}{\frac{1}{4}}=1\]
 \[\frac{800+x}{1000}\cdot\frac{23}{36}=1\]
 \[800+x=\frac{36000}{23}\]
 \[x=\frac{17600}{23}.\]

 We can't have a part of a worker, so we take the ceiling of $x$, which we find to be $\boxed{766}$.

 -flyhawkeye

 





\textbf{Solution 3 (Very fast, the one I used on the test)}

Suppose $1000$ workers can complete one quarter of the job in one day. After the first day, there were $900$ workers remaining so the second quarter was completed in $\frac{10}{9}$ days. Now there are only $800$ workers remaining so the third quarter can be completed in $\frac{10}{8}$ days. It has been $1+\frac{10}{9}+\frac{10}{8}=\frac{121}{36}$ days since the job started, and we still have $\frac{23}{36}$ days to complete the last quarter of the job, so we need $\frac{36}{23}$ of $1000$ workers, or $1565.2173913$. Since we can’t hire fractional workers, $1566-800=\boxed{766}$ must be the answer.

 




	\item (\textit{AIME II 2004}) Three clever monkeys divide a pile of bananas. The first monkey takes some bananas from the pile, keeps three-fourths of them, and divides the rest equally between the other two. The second monkey takes some bananas from the pile, keeps one-fourth of them, and divides the rest equally between the other two. The third monkey takes the remaining bananas from the pile, keeps one-twelfth of them, and divides the rest equally between the other two. Given that each monkey receives a whole number of bananas whenever the bananas are divided, and the numbers of bananas the first, second, and third monkeys have at the end of the process are in the ratio $3: 2: 1,$ what is the least possible total for the number of bananas?


 
 \textbf{Answer: }408



\textbf{Solution 1}

In this solution, you start out the same as solution one. Convert everything into the fractions of largest denominator terms (this is necessary) until you get 

 $27x=11z$$27y=13z$

 While solving, and make sure to leave a list of numbers that must divide $x$, $y$, and $z$. For example, just by looking at the basic fractions you receive from writing the starting equations, 24 divides $z$, 8 divides $y$ and $x$. In the expression above, it's also clear that 27 divides $z$, 13 divides $y$, and 11 divides $z$. You might be wondering why I wrote the expression in terms of z. That's because $z$ has the largest divisor. The LCM of all it's divisors shows that $z$ must be divisible by 216. The total amount of bananas can be found to equal to $17z/9$. This means there are two possible solutions under 1000: 408 and 816. Trial and error can be done quickly to find the smallest possible solution to be $\boxed{408}$.

 
-jackshi2006

 





\textbf{Solution 2 (A lot of Thinking Required dEfInItELy )}

Let $A,B,C$ be the fraction of bananas taken by the first, second, and third monkeys respectively. Then we have the system of equations 
 \[\frac{3}{4}A+\frac{3}{8}B+\frac{11}{24}C=\frac{1}{2}\] 
 \[\frac{1}{8}A+\frac{1}{4}B+\frac{11}{24}C=\frac{1}{3}\] 
 \[\frac{1}{8}A+\frac{3}{8}B+\frac{2}{24}C=\frac{1}{6}.\] Solve this your favorite way to get that
 \[(A,B,C)=\left( \frac{11}{51}, \frac{13}{51}, \frac{9}{17} \right).\] We need the amount taken by the first and second monkeys to be divisible by 8 and the third by 24 (but for the third, we already have divisibility by 9). Thus our minimum is $8 \cdot 51 = \boxed{408}.$

 




	\item (\textit{AIME II 2004}) $ABCD$ is a rectangular sheet of paper that has been folded so that corner $B$ is matched with point $B'$ on edge $AD.$ The crease is $EF,$ where $E$ is on $AB$ and $F$ is on $CD.$ The dimensions $AE=8, BE=17,$ and $CF=3$ are given. The perimeter of rectangle $ABCD$ is $m/n,$ where $m$ and $n$ are relatively prime positive integers. Find $m+n.$


 
                    \begin{figure}[H]
                        \centering
                        \includegraphics[width=0.5\textwidth]{Problem/AIME_II_2004_7_Problem_Figure_1.png}
                    \end{figure}
                    


 
 \textbf{Answer: }293



\textbf{Solution 1 (Synthetic)}


                    \begin{figure}[H]
                        \centering
                        \includegraphics[width=0.5\textwidth]{Solution/AIME_II_2004_7_Solution_Figure_1.png}
                    \end{figure}
                    Since $EF$ is the perpendicular bisector of $\overline{BB'}$, it follows that $BE = B'E$ (by SAS). By the \href{https://artofproblemsolving.com/wiki/index.php?title=Pythagorean\_Theorem}{Pythagorean Theorem}, we have $AB' = 15$. Similarly, from $BF = B'F$, we have
 \begin{align*} BC^2 + CF^2 = B'D^2 + DF^2 &\Longrightarrow BC^2 + 9 = (BC - 15)^2 + 484 \\ BC  &= \frac{70}{3} \end{align*}
Thus the perimeter of $ABCD$ is $2\left(25 + \frac{70}{3}\right) = \frac{290}{3}$, and our answer is $m+n=\boxed{293}$.

 





\textbf{Solution 2 (analytic)}

Let $A = (0,0), B=(0,25)$, so $E = (0,8)$ and $F = (l,22)$, and let $l = AD$ be the length of the rectangle. The \href{https://artofproblemsolving.com/wiki/index.php?title=Slope}{slope} of $EF$ is $\frac{14}{l}$ and so the equation of $EF$ is $y -8 = \frac{14}{l}x$. We know that $EF$ is perpendicular to and bisects $BB'$. The slope of $BB'$ is thus $\frac{-l}{14}$, and so the equation of $BB'$ is $y -25 = \frac{-l}{14}x$. Let the point of intersection of $EF, BB'$ be $G$. Then the y-coordinate of $G$ is $\frac{25}{2}$, so
 \begin{align*} \frac{14}{l}x &= y-8 = \frac{9}{2}\\ \frac{-l}{14}x &= y-25 = -\frac{25}{2}\\ \end{align*}
Dividing the two equations yields

 $l^2 = \frac{25 \cdot 14^2}{9} \Longrightarrow l = \frac{70}{3}$The answer is $\boxed{293}$ as above.

 





\textbf{Solution 3 (Coordinate Bashing)}

Firstly, observe that if we are given that $AE=8$ and $BE=17$, the length of the triangle is given and the height depends solely on the length of $CF$. Let Point $A = (0,0)$. Since $AE=8$, point E is at (8,0). Next, point $B$ is at $(25,0)$ since $BE=17$ and point $B'$ is at $(0,15)$ since $BE=AE$ by symmetry. Draw line segment $BB'$. Notice that this is perpendicular to $EF$ by symmetry. Next, find the slope of EB, which is $\frac{15}{25}=\frac{3}{5}$. Then, the slope of $EF$ is -$\frac{5}{3}$. 

 Line EF can be written as y=$-\frac{5}{3}x+b$. Plug in the point $(8,0)$, and we get the equation of EF to be y=$_\frac{5}{3}x+\frac{40}{3}$. Since the length of $AB$=25, a point on line $EF$ lies on $DC$ when $x=25-3=22$. Plug in $x=22$ into our equation to get $y=-\frac{70}{3}$. $|y|=BC=\frac{70}{3}$. Therefore, our answer is $2(AB+BC)=2\left(25+\frac{70}{3}\right)=2\left(\frac{145}{3}\right)=\frac{290}{3}= \boxed{293}$.

 





\textbf{Solution 4 (Trig)}

Firstly, note that $B'E=BE=17$, so $AB'=\sqrt{17^2-8^2}=15$. Then let $\angle BEF=\angle B'EF=\theta$, so $\angle B'EA = \pi-2\theta$. Then $\tan(\pi-2\theta)=\frac{15}{8}$, or 

 
 \[\frac{2\tan(\theta)}{\tan^2(\theta)-1}=\frac{15}{8}\] using supplementary and double angle identities. Multiplying though and factoring yields 

 
 \[(3\tan(\theta)-5)(5\tan(\theta)+3)=0\]

 It is clear from the problem setup that $0<\theta<\frac\pi2$, so the correct value is $\tan(\theta)=\frac53$. Next, extend rays $\overrightarrow{BC}$ and $\overrightarrow{EF}$ to intersect at $C'$. Then $\tan(\theta)=\frac{BC'}{17}=\frac53$, so $BC'=\frac{85}{3}$. By similar triangles, $CC'=\frac{3}{17}BC'=\frac{15}{3}$, so $BC=\frac{70}{3}$. The perimeter is $\frac{140}{3}+50=\frac{290}{3}\Longrightarrow \boxed{293}$

 An even faster way to finish is, to draw a line segment $FF'$ where $F'$ is a point on $EB$ such that $FF'$ is perpendicular to $EB$. This makes right triangle $FF'E$, Also, note that $F'B$ has length of $3$ (draw the diagram out, and note the $F'B =FC$). From here, through $\tan \theta = \frac{5}{3}$, we can note that $\frac{FF'}{EF'} = \frac{5}{3} \implies \frac{FF'}{14} = \frac{5}{3} \implies FF' = \frac{70}{3}$. $FF'$ is parallel and congrurent to $CB$ and $AD$, and hence we can use this to calculate the perimeter. The perimeter is simply $\frac{70}{3} + \frac{70}{3} + 25 + 25 = \frac{290}{3} \Longrightarrow \boxed{293}$

 





\textbf{Solution 5 (Fast, Pythagorean)}

Use the prepared diagram for this solution.

 Call the intersection of DF and B'C' G. AB'E is an 8-15-17 right triangle, and so are B'DG and C'FG. Since C'F is 3, then using the properties of similar triangles GF is 51/8. DF is 22, so DG is 125/8. Finally, DB can to calculated to be 25/3.Add all the sides together to get $\boxed{293}$.

 -jackshi2006

 





\textbf{Solution 6(fast as wind[rufeng])}

Call the intersection of $B'C'$, $BC$, and $EF$ $G$. Since $FCBE$ and $FC'B'E$ are congruent, we know that the three lines intersect.We already know $AB$ so we just need to find $CB$, call it $x$. Drop an altitude from $F$ to $AB$ and call it $H$. $EH=EB-FC=14$. Using Pythagorean Theorem, we have $EF=\sqrt{x^2+14^2}$. Triangles $EFH$ and $EGB$ are similar (AA), so we get
 \[\frac{HF}{BG}=\frac{EH}{EB}\]
 \[\frac{x}{x+GC}=\frac{14}{17}\]Simplify and we get $GC=\frac{3x}{14}$. 

 We find the area of $FCBE$ by using the fact that it is a trapezoid. $[FCBE]=\frac{(3+17)x}{2}=10x$

 A different way to find the area:$[FCBE]=\frac{1}{2} EG\cdot($height of $EGB$ with $EG$ as base$)-[FGC]$

 Since $GBE$ and $G'B'E$ are congruent(SAS), their height from $EG$ is the same. $B'B=\sqrt{AB'^2+AB^2}=5\sqrt{34}$.$EG=\sqrt{EB^2+BG^2}=\sqrt{(\frac{17x}{14})^2+17^2}=17\sqrt{\frac{x^2}{196}+1}$

 
 \[[FCBE]=\frac{1}{2} \cdot 17 \cdot \sqrt{\frac{x^2}{196}+1} \cdot \frac{5\sqrt{34}}{2}-\frac{9x}{28}\]
 \[280x+9x=7\cdot 5 \cdot \sqrt{34} \cdot 17 \cdot \sqrt{\frac{x^2}{196}+1}\]
 \[17^4 x^2=49 \cdot 25 \cdot 34 \cdot 17^2 \cdot (\frac{x^2}{196}+1)\]
 \[17x^2=\frac{25}{2}x^2+2450\]
 \[x=\frac{70}{3}\]

 The perimeter is $\frac{140}{3}+50=\frac{290}{3},$ so our answer is $\boxed{293}$.

 





\textbf{Solution 7 (Similar to solution 5, more in depth)}

Let the endpoint of the intersection of the fold near $F$ be $G$. Since trapezoid $BCFE$ is folded, it is congruent to trapezoid $B'C'FE$. Therefore, $BE=B'E=17$. Since $\triangle AB'E$ is a right triangle, $AB'=15$ from the pythagorean theorem. From here, we can see that triangles $\triangle AEB \sim \triangle DGB' \sim \triangle C'GF$ by AA similarity. From here, we find $BC$ from a lot of similarities. Let $BC=x$.

 Since $\triangle ABE' \sim \triangle DGB'$:

 
 \[\frac {AE}{AB'} = \frac{DB}{DG}\]

 
 \[\frac {8}{15} = \frac {x-15}{DG}\]

 
 \[DG = \frac {15(x-15)}{8}\]

 
 \[GF = DC-DG-FC\]

 
 \[GF = \frac{-15x+401}{8}\]

 Since $\triangle ABE' \sim \triangle C'GF'$,

 
 \[\frac {AE}{B'E} = \frac {C'F}{GF}\]

 
 \[\frac {8}{17} = \frac{3}{\frac {-15x+401}{8}}\]

 from which we get $x= \frac {70}{3}$.

 Finally, our answer is $2(\frac {70}{3}) + 2(25)=\frac {290}{3}$, which is $290+3=\boxed{293}$.

 




	\item (\textit{AIME II 2004}) How many positive integer divisors of $2004^{2004}$ are divisible by exactly 2004 positive integers?


 
 \textbf{Answer: }054



\textbf{Solution 1}

The \href{https://artofproblemsolving.com/wiki/index.php?title=Prime\_factorization}{prime factorization} of 2004 is $2^2\cdot 3\cdot 167$.  Thus the prime factorization of $2004^{2004}$ is $2^{4008}\cdot 3^{2004}\cdot 167^{2004}$.

 We can \href{https://artofproblemsolving.com/wiki/index.php?title=Divisor\_function}{ count the number of divisors} of a number by multiplying together one more than each of the \href{https://artofproblemsolving.com/wiki/index.php?title=Exponentiation}{ exponents} of the prime factors in its prime factorization.  For example, the number of divisors of $2004=2^2\cdot 3^1\cdot 167^1$ is $(2+1)(1+1)(1+1)=12$.  

 A positive integer divisor of $2004^{2004}$ will be of the form $2^a\cdot 3^b\cdot 167^c$.  Thus we need to find how many $(a,b,c)$ satisfy 

 $(a+1)(b+1)(c+1)=2^2\cdot 3\cdot 167.$We can think of this as \href{https://artofproblemsolving.com/wiki/index.php?title=Partition}{partitioning} the exponents to $a+1,$ $b+1,$ and  $c+1$.  So let's partition the 2's first.  There are two 2's so this is equivalent to partitioning two items in three containers.  We can do this in ${4 \choose 2} = 6$ ways.  We can partition the 3 in three ways and likewise we can partition the 167 in three ways.  So we have $6\cdot 3\cdot 3 = \boxed{54}$ as our answer.

 





\textbf{Solution 2 (bash)}

Clearly we need to find a group of numbers that multiply to 2004. We can list them all out since we know that 2004 is only $167 \cdot 2^2 \cdot 3$.

 167, 2, 2, 3

 4, 3, 167

 12, 167

 4, 501

 2, 1002

 2, 3, 334

 2, 2, 501*

 6, 2, 167

 3, 668

 6, 334

 2004*

 To begin, the first multiple doesn't work because there are only 3 prime divisors of 2004. We can apply all multiples because the prime factorization of $2004^{2004}$ is $2^{4008} \cdot 3^{2004} \cdot 167^{2004}.$ Every multiple has six ways of distributing numbers to become powers of 167, 3, and 2, except for the ones with a star. For a single power of 2004, we have three choices (2, 3, and 167) to give a power of 2003 to.For 2, 2, 501, there are three choices to give a power of 500 to and the rest get a power of 1.

 Therefore we have $8 \cdot 6$ normal multiples and $3 \cdot 2$ "half" multiples. Sum to get $\boxed{54}$ as our answer.

 


 





\textbf{Solution 3}

we want $167^a \cdot 3^b \cdot 2^c$ such that this number has exactly $2004$ positive factors. Note that $2004$ has $(1 + 1) \cdot (1 + 1) \cdot (2 + 1) = 12$ factors.

 


 \textbf{Case 1:} $(a + 1)(b + 1)(c + 1) = 2004$, and no variable equals $0.$

 


 The three variables can either be the three elements of $\{166, 1, 5\}, \{500, 1, 1\}, \{333, 2, 1\},$ or $\{166, 3, 2\}.$ They can be permuted in $6 + 3 + 6 + 6 = 21$ ways.

 


 \textbf{Case 2: } $(x + 1)(y + 1) = 2004$, where $x$ and $y$ are two variables chosen out of $\{a, b, c\}$ and none of them equal $0$

 


 There are $12 - 2 = 10$ ways to choose the values of $x$ and $y$. There are $\binom{3}{2} = 3$ ways to choose the variables $x$ and $y$ out of  $\{a, b, c\}$, giving $30$ ways.

 


 \textbf{Case 3:} $(x + 1) = 2004$

 


 $x$ must be $2003$. We can let $x = a,b,$ or $c$, so this case gives us $3$ ways.

 
The answer is $21 + 30 + 3 = \boxed{54}.$

 





\textbf{Solution 4 (Stars and Bars)}


 \[2004^{2004}\]
 \[= (2^2 \cdot 3 \cdot 167)^{2004}\]
 \[= 2^{4008} \cdot 3^{2004} \cdot 167^{2004}\]
 \[2^{i}3^{j}167^{k} \ | \ 2004^{2004}\]
 \[0 \le i \le 4008, 0 \le j \le 2004, 0 \le k \le 2004\]
 \[i,j,k \not \ge 2004 \ \text{otherwise} \ (i+1) > 2004\]
 \[0 \le i, j, k < 2004\]
 \[(i+1)(j+1)(k+1) = 2004\]

 Each of these must contain values that multiply to $2^2, 3^1, 167^1$. Consider the exponents of each of the values $i+1, j+1, k+1$. Then, we see that if each of $i+1 = a, j+1 = b, k+1 = c$ exists as a prime factorization $a,b,c = 2^{x_1,x_4,x_7}3^{x_2,x_5,x_8}167^{x_3,x_6,x_9}$. Therefore, we see that the solutions for $a$ exist as $x_1 + x_4 + x_7 = 2$, for $b$, $x_2 + x_5 + x_8 = 1$, and for $c$, $x_3 + x_6 + x_9 = 1$. Through stars and bars, we have $a$ has 6 choices, $b$ has 3, and so does $c$. Then, $abc$ has $6 \cdot 3 \cdot 3 = \boxed{054}$ total cases.

 




	\item (\textit{AIME II 2004}) A sequence of positive integers with $a_1=1$ and $a_9+a_{10}=646$ is formed so that the first three terms are in geometric progression, the second, third, and fourth terms are in arithmetic progression, and, in general, for all $n\ge1,$ the terms $a_{2n-1}, a_{2n}, a_{2n+1}$ are in geometric progression, and the terms $a_{2n}, a_{2n+1},$ and $a_{2n+2}$ are in arithmetic progression. Let $a_n$ be the greatest term in this sequence that is less than 1000. Find $n+a_n.$


 
 \textbf{Answer: }973



\textbf{Solution 1}

Let $x = a_2$; then solving for the next several terms, we find that $a_3 = x^2,\ a_4 = x(2x-1),\ a_5$ $= (2x-1)^2,\ a_6$ $= (2x-1)(3x-2)$, and in general, $a_{2n} = f(n-1)f(n)$, $a_{2n+1} = f(n)^2$, where $f(n) = nx - (n-1)$.<sup>[\href{https://artofproblemsolving.com\#note\_1}{1}]</sup> 

 From 
 \[a_9 + a_{10} = f(4)^2 + f(4)f(5) = (4x-3)(9x-7) = 646 = 2\cdot 17 \cdot 19\], we find that by either the \href{https://artofproblemsolving.com/wiki/index.php?title=Quadratic\_formula}{quadratic formula} or trial-and-error/modular arithmetic that $x=5$. Thus $f(n) = 4n+1$, and we need to find the largest $n$ such that either $f(n)^2\, \mathrm{or}\, f(n)f(n-1) < 1000$. This happens with $f(7)f(8) = 29 \cdot 33 = 957$, and this is the $2(8) = 16$th term of the sequence.

 The answer is $957 + 16 = \boxed{973}$.

 \href{https://artofproblemsolving.com\#ref\_1}{\textasciicircum{}} We can show this by simultaneous \href{https://artofproblemsolving.com/wiki/index.php?title=Induction}{induction}: since 
 \begin{align*}a_{2n} &= 2a_{2n-1} - a_{2n-2} = 2a_{2(n-1)+1} - a_{2(n-1)} \\ &= 2f(n-1)^2 - f(n-2)f(n-1) \\ &= f(n-1)[2f(n-1) - f(n-2)] \\ &= f(n-1)[(2n-2-n+2)x-(2n-4-n+3)] \\ &= f(n-1)f(n) \end{align*}
and
 \begin{align*}a_{2n+1} &= \frac{a_{2n}^2}{a_{2n-1}} = \frac{f(n-1)^2f(n)^2}{f(n-1)^2} = f(n)^2 \\ \end{align*}


 





\textbf{Solution 2}

Let $x = a_2$. It is apparent that the sequence grows relatively fast, so we start trying positive integers to see what $x$ can be. Finding that $x = 5$ works, after bashing out the rest of the terms we find that $a_{16} = 957$ and $a_{17} = 1089$, hence our answer is $957 + 16 = \boxed{973}$.

 





\textbf{Solution 3}

We can find the value of $a_{9}$ by its bounds using three conditions:

 
            \begin{enumerate}
                \item $0<a_{8} = 2a_{9}-a_{10}$\item $a_{10} < a_{11}$ (note that the sequence must be increasing on all terms, not monotonically increasing) $a_{10} < \frac{a_{10}^2}{a_{9}} \rightarrow a_{9} < a_{10}$\item $a_{11} = \frac{a_{10}^2}{a_{9}} = \frac{(646-a_{9})^2}{a_{9}}$, so necessarily $a_{9}$ is a factor of $646^2$, which factorizes to $2^2\cdot 17^2 \cdot 19^2$
            \end{enumerate}
            Rearranging conditions 1 and 2, we get:

 
 \[\frac{646}{3} < a_{9} < \frac{646}{2}\]

 trying all the terms from the third condition, it is clear that $a_9 = 289$ is the only solution. Then we can calculate the next few terms from there since we have $a_{10}$ as well, to find that $a_{16} = 957$ and $a_{17} = 1089$, thus we have our answer of $957 + 16 = \boxed{973}$.

 





\textbf{Solution 4 (quadratic)}

Let $a_2 = x$. We will write the first 10 terms in terms of $x$ (this will require some rigorous polynomial division):

 




	\item (\textit{AIME II 2004}) Let $S$ be the set of integers between 1 and $2^{40}$ whose binary expansions have exactly two 1's. If a number is chosen at random from $S,$ the probability that it is divisible by 9 is $p/q,$ where $p$ and $q$ are relatively prime positive integers. Find $p+q.$


 
 \textbf{Answer: }913



\textbf{Solution 1 (Rigorous)}

Any number from 1 to $2^{40}$ can be represented in binary with 40 digits (because $2^{40}$ has 41) with leading zeroes if necessary. Therefore the number of sets where there are exactly two 1’s in this binary representation is just $\binom {40}{2}$ because we’re choosing 2 1s to go in 40 digit slots. This is equal to 780; we have found $q$, our denominator.

 With these particular numbers, when converting back to decimal note that the 38 zeroes all cancel out. Let the first 1 when reading the number in binary from left to right be in position $y$, I.e. when converting to decimal, it will have value $2^{y}$. Similarly, let the second 1 have position $x$ (note $x<y$); then the decimal value of that 1 will be $2^{x}$. For example, with the binary string 0001001000 $y$ is 6 and $x$ is 3 (note that it is zero indexed).

 It should also be noted that our constraints on $x$ and $y$ are $0 \leq x < y < 40$ and $x, y \in \mathbb{Z}$.

 The numbers that satisfy the constraints are therefore equal to $2^{x}+2^{y} = 2^{x} \cdot (1+2^{y-x})$. The conditions say that this is divisible by 9. $2^{x}$ is clearly never divisible by 9, so we can just focus on $1+2^{y-x}$. 
 \[1+2^{y-x} \equiv 0 \pmod 9 \implies 2^{y-x} \equiv -1 \pmod 9 \implies 2^{y-x} \equiv 8 \pmod 9\]

 $y-x = 3$ clearly satisfies the congruence, but are there any more? Well, we see that any time we raise -1 (in the second congruence above) to an odd power we get -1 again, so if we raise $2^{3}$, which we know already works, to an odd power, we will also satisfy the congruence. Thus, $2^{3}, 2^{9}, 2^{15},$ etc. work; or, 
 \[y-x \equiv 3 \pmod 6\]

 To count the numbers that satisfy the condition, we need to find the amount of numbers that satisfy $2^{x} \cdot (1+2^{y-x}) \leq 2^{40}$, and $y-x \equiv 3 \pmod 6 = 6n +3$ for some whole number $n$. 

 Since we know $1+2^{y-x} = 1+2^{6n+3}$ is positive, we can divide both sides of the inequality by it;
 \[2^{x} \cdot (1+2^{y-x}) \leq 2^{40} \implies 2^{x} \leq \frac{2^{40}}{1+2^{y-x}} \implies 2^{x} \leq \frac{2^{40}}{1+2^{6n+3}}\]
 \[\implies 2^{x} < \frac{2^{40}}{2^{6n+3}} \implies 2^{x} < 2^{40-(6n+3)}\implies 2^{x} < 2^{37-6n} \implies x < 37 - 6n\]

 Finally, the number of possible values of $x$ is $37-6n$ because $x$ can also be 0 and must be an integer.Therefore, as we increase $n$ from zero upwards, we get the total number of possibilities for $x$ is $37+31+25+\ldots+1 = 7 + 36 + 30 + 24 + \ldots + 6 + 0 = 7 + 6 \cdot (6 + 5 + 4\ldots + 1)$

 $= 7 + 6 \cdot 3 \cdot 7 = 7 + 18 \cdot 7 = 19 \cdot 7 = 133$.

 Since 133 and 780 are relatively prime, the answer is $\frac{133}{780}$ and our final answer is $133+780 = \boxed{913}$

 





\textbf{Solution 2}

A positive integer $n$ has exactly two 1s in its binary representation exactly when $n = 2^j + 2^k$ for $j \neq k$ \href{https://artofproblemsolving.com/wiki/index.php?title=Nonnegative}{nonnegative} integers.  Thus, the \href{https://artofproblemsolving.com/wiki/index.php?title=Set}{set} $S$ is equal to the set $\{n \in \mathbb{Z} \mid n = 2^j + 2^k \,\mathrm{ and }\, 0 \leq j < k \leq 39\}$.  (The second condition ensures simultaneously that $j \neq k$ and that each such number less than $2^{40}$ is counted exactly once.)  This means there are ${40 \choose 2} = 780$ total such numbers.

 Now, consider the powers of $2$ \href{https://artofproblemsolving.com/wiki/index.php?title=Modular\_arithmetic}{ mod} $9$: $2^{6n} \equiv 1, 2^{6n + 1} \equiv 2, 2^{6n + 2} \equiv 4, 2^{6n + 3} \equiv 8 \equiv -1,$ $2^{6n + 4} \equiv 7 \equiv -2,$ $2^{6n + 5} \equiv 5 \equiv -4 \pmod 9$.  

 It's clear what the pairs $j, k$ can look like.  If one is of the form $6n$ (7 choices), the other must be of the form $6n + 3$ (7 choices).  If one is of the form $6n + 1$ (7 choices) the other must be of the form $6n + 4$ (6 choices).  And if one is of the form $6n + 2$ (7 choices), the other must be of the form $6n + 5$ (6 choices).  This means that there are $7\cdot 7 + 7\cdot 6 + 7\cdot 6 = 49 + 42 +42 = 133$ total "good" numbers.

 The probability is $\frac{133}{780}$, and the answer is $133 + 780 = \boxed{913}$.

 -minor edit by Mathkiddie

 





\textbf{Solution 3 (similar to 2)}

As mentioned above, there are 780 possible combinations. Since we are essentially adding two powers of two together, thinking about the properties of this sum organizes our solution. All powers of two are even except for $2^0$, so we begin with labeling an entire group "where one of the 1s is in the rightmost spot". In an equation, $2^n+1$ is congruent to 0 mod 9. Instantly we think of $2^3=8$, and trial and error soon gives us $2^9=512$. Notice a pattern? Trying 2\textasciicircum{}15 = 32768 also works. You could go on, but basically all powers of two 3 mod 6 are 1 mod 9. (This is proven with the fact that $2^6$ is 1 mod 9) There are seven 3 mod 6 powers of 2 from 1-39, so for our first category there are seven "good" combinations.

 
Next, we look at powers of two exclusive from 1-39. Expressed in an equation, $2^n+2^m$ is congruent to 0 mod 9. Since nine is a relatively small number, we can treat powers of two as remainders of nine, to find "good" combinations quickly. Making a list of powers of two from 1-6 shows that the remainders when divided by nine are 2, 4, 8, 7, 5, and 1. This repeats itself every 6 powers. With simple counting we see there are seven powers of two that give remainders of 2, 4, and 8 each. Also there are six powers of two two that give remainders of 7, 5, and 1 each. Notice that each trio of powers correspond with another in the opposite trio. WLOG, if you choose two, (which there are seven to select) there are six different sevens to pair. This gives $7*6$. Repeat for the other two pairs.

 In the end there are $7+42+42+42$ "good" combinations, for an answer of $133 + 780 = \boxed{913}$.

 
-jackshi2006

 





\textbf{Solution 4}

Notice that in any binary expression, when we take it modulo $9$ and look at it in groups of 3 starting from the right, we get bases with modulo's $4,2,1$ in the first group and $-4,-2,-1$ in the second group. For our number to be divisible by $9$, we have to get one of the positive multipliers and its opposite to cancel out. There are a total of $7$ complete positive groups and $6$ complete negative groups and then $2^{39}$ is left which is $-1$. If we select a $1$, we have $7$ ways to negate it and $7$ ways to select it giving $49$. If we select a $2$, we have $6$ ways to negate it and $7$ ways to select it giving $42$. The same follows for selecting a $4$. In total, we need to choose $2$ out of the total $40$ digits to "activate" and turn into a $1$. In total we get 
 \[\frac{49+42+42}{{40\choose{2}}} = \frac{133}{780}\] Adding, we get $133 + 780 = \boxed{913}$.

 -Vedoral

 




	\item (\textit{AIME II 2004}) A right circular cone has a base with radius 600 and height $200\sqrt{7}.$ A fly starts at a point on the surface of the cone whose distance from the vertex of the cone is 125, and crawls along the surface of the cone to a point on the exact opposite side of the cone whose distance from the vertex is $375\sqrt{2}.$ Find the least distance that the fly could have crawled.


 
 \textbf{Answer: }625



\textbf{Solution 1}

The easiest way is to unwrap the cone into a circular sector. Center the sector at the origin with one radius on the positive $x$-axis and the angle $\theta$ going counterclockwise. The circumference of the base is $C=1200\pi$. The sector's radius (cone's sweep) is $R=\sqrt{r^2+h^2}=\sqrt{600^2+(200\sqrt{7})^2}=\sqrt{360000+280000}=\sqrt{640000}=800$. Setting $\theta R=C\implies 800\theta=1200\pi\implies\theta=\frac{3\pi}{2}$.

 If the starting point $A$ is on the positive $x$-axis at $(125,0)$ then we can take the end point $B$ on $\theta$'s bisector at $\frac{3\pi}{4}$ radians along the $y=-x$ line in the second quadrant. Using the distance from the vertex puts $B$ at $(-375,-375)$. Thus the shortest distance for the fly to travel is along segment $AB$ in the sector, which gives a distance $\sqrt{(-375-125)^2+(-375-0)^2}=125\sqrt{4^2+3^2}=\boxed{625}$.

 





\textbf{Solution 2}


                    \begin{figure}[H]
                        \centering
                        \includegraphics[width=0.5\textwidth]{Solution/AIME_II_2004_11_Solution_Figure_1.png}
                    \end{figure}
                    To find the shortest length from the red to blue points, the net of the side of the cone could be drawn.

 
                    \begin{figure}[H]
                        \centering
                        \includegraphics[width=0.5\textwidth]{Solution/AIME_II_2004_11_Solution_Figure_2.png}
                    \end{figure}
                    The angle $YVX$ is equal to $360^\circ \cdot \frac{1200\pi}{1600\pi} \cdot \frac{1}{2}$, or $135^\circ$. Therefore, the law of cosines could be utilized.
 \[AB = \sqrt{(375\sqrt{2})^2 + 125^2 - 2 \cdot (375\sqrt{2})(125)(\cos 135^\circ)} = \boxed{625}\]

 




	\item (\textit{AIME II 2004}) Let $ABCD$ be an isosceles trapezoid, whose dimensions are $AB = 6, BC=5=DA,$and $CD=4.$ Draw circles of radius 3 centered at $A$ and $B$ and circles of radius 2 centered at $C$ and $D.$ A circle contained within the trapezoid is tangent to all four of these circles. Its radius is $\frac{-k+m\sqrt{n}}p,$ where $k, m, n,$ and $p$ are positive integers, $n$ is not divisible by the square of any prime, and $k$ and $p$ are relatively prime. Find $k+m+n+p.$


 
 \textbf{Answer: }134



\textbf{Solution}

Let the radius of the center circle be $r$ and its center be denoted as $O$.

 
                    \begin{figure}[H]
                        \centering
                        \includegraphics[width=0.5\textwidth]{Solution/AIME_II_2004_12_Solution_Figure_1.png}
                    \end{figure}
                    Clearly line $AO$ passes through the point of tangency of circle $A$ and circle $O$. Let $y$ be the height from the base of the trapezoid to $O$. From the \href{https://artofproblemsolving.com/wiki/index.php?title=Pythagorean\_Theorem}{Pythagorean Theorem}, 
 \[3^2 + y^2 = (r + 3)^2 \Longrightarrow y = \sqrt {r^2 + 6r}.\]

 We use a similar argument with the line $DO$, and find the height from the top of the trapezoid to $O$, $z$, to be $z = \sqrt {r^2 + 4r}$.

 Now $y + z$ is simply the height of the trapezoid. Let $D'$ be the foot of the \href{https://artofproblemsolving.com/wiki/index.php?title=Perpendicular}{perpendicular} from $D$ to $AB$; then $AD' = 3 - 2 = 1$. By the Pythagorean Theorem, $(AD')^2 + (DD')^2 = (AD)^2 \Longrightarrow DD' = \sqrt{24}$ so we need to solve the equation $\sqrt {r^2 + 4r} + \sqrt {r^2 + 6r} = \sqrt {24}$. We can solve this by moving one radical to the other side, and squaring the equation twice to end with a \href{https://artofproblemsolving.com/wiki/index.php?title=Quadratic\_equation}{quadratic equation}. 

 Solving this, we get $r = \frac { - 60 + 48\sqrt {3}}{23}$, and the answer is $k + m + n + p = 60 + 48 + 3 + 23 = \boxed{134}$.

 




	\item (\textit{AIME II 2004}) Let $ABCDE$ be a convex pentagon with $AB || CE, BC || AD, AC || DE, \angle ABC=120^\circ, AB=3, BC=5,$ and $DE = 15.$ Given that the ratio between the area of triangle $ABC$ and the area of triangle $EBD$ is $m/n,$ where $m$ and $n$ are relatively prime positive integers, find $m+n.$


 
 \textbf{Answer: }484



\textbf{Solution}

Let the intersection of $\overline{AD}$ and $\overline{CE}$ be $F$. Since $AB \parallel CE, BC \parallel AD,$ it follows that $ABCF$ is a \href{https://artofproblemsolving.com/wiki/index.php?title=Parallelogram}{parallelogram}, and so $\triangle ABC \cong \triangle CFA$. Also, as $AC \parallel DE$, it follows that $\triangle ABC \sim \triangle EFD$. 

 
                    \begin{figure}[H]
                        \centering
                        \includegraphics[width=0.5\textwidth]{Solution/AIME_II_2004_13_Solution_Figure_1.png}
                    \end{figure}
                    By the \href{https://artofproblemsolving.com/wiki/index.php?title=Law\_of\_Cosines}{Law of Cosines}, $AC^2 = 3^2 + 5^2 - 2 \cdot 3 \cdot 5 \cos 120^{\circ} = 49 \Longrightarrow AC = 7$. Thus the length similarity ratio between $\triangle ABC$ and $\triangle EFD$ is $\frac{AC}{ED} = \frac{7}{15}$. 

 Let $h_{ABC}$ and $h_{BDE}$ be the lengths of the \href{https://artofproblemsolving.com/wiki/index.php?title=Altitude}{altitudes} in $\triangle ABC, \triangle BDE$ to $AC, DE$ respectively. Then, the ratio of the areas $\frac{[ABC]}{[BDE]} = \frac{\frac 12 \cdot h_{ABC} \cdot AC}{\frac 12 \cdot h_{BDE} \cdot DE} = \frac{7}{15} \cdot \frac{h_{ABC}}{h_{BDE}}$. 

 However, $h_{BDE} = h_{ABC} + h_{CAF} + h_{EFD}$, with all three heights oriented in the same direction. Since $\triangle ABC \cong \triangle CFA$, it follows that $h_{ABC} = h_{CAF}$, and from the similarity ratio, $h_{EFD} = \frac{15}{7}h_{ABC}$. Hence $\frac{h_{ABC}}{h_{BDE}} = \frac{h_{ABC}}{2h_{ABC} + \frac {15}7h_{ABC}} = \frac{7}{29}$, and the ratio of the areas is $\frac{7}{15} \cdot \frac 7{29} = \frac{49}{435}$. The answer is $m+n = \boxed{484}$.

 \textbf{Additional Trigonometry-Free Alternative} Instead of using the Law of Cosines, we can draw a line perpendicular to line BC down from point A until it intersects BC at a point $P$. Since $\angle PBA = 60^{\circ}$, we can use the $30-60-90$ triangle to deduce that $PB = \frac{3}{2}$, and $PA = \frac{3\sqrt{3}}{2}$. From here, we can use Pythagorean theorem to deduce that $AC = 7$. Then, we can follow with the rest of the solution above.

 




	\item (\textit{AIME II 2004}) Consider a string of $n$ $7$'s, $7777\cdots77,$ into which $+$ signs are inserted to produce an arithmetic expression. For example, $7+77+777+7+7=875$ could be obtained from eight $7$'s in this way. For how many values of $n$ is it possible to insert $+$ signs so that the resulting expression has value $7000$?


 
 \textbf{Answer: }108



\textbf{Solution 1}

Suppose we require $a$ $7$s, $b$ $77$s, and $c$ $777$s to sum up to $7000$ ($a,b,c \ge 0$). Then $7a + 77b + 777c = 7000$, or dividing by $7$, $a + 11b + 111c = 1000$. Then the question is asking for the number of values of $n = a + 2b + 3c$.

 Manipulating our equation, we have $a + 2b + 3c = n = 1000 - 9(b + 12c) \Longrightarrow 0 \le 9(b+12c) < 1000$. Thus the number of potential values of $n$ is the number of multiples of $9$ from $0$ to $1000$, or $112$.

 
However, we forgot to consider the condition that $a \ge 0$. For a solution set $(b,c): n=1000-9(b+12c)$, it is possible that $a = n-2b-3c < 0$ (for example, suppose we counted the solution set $(b,c) = (1,9) \Longrightarrow n = 19$, but substituting into our original equation we find that $a = -10$, so it is invalid). In particular, this invalidates the values of $n$ for which their only expressions in terms of $(b,c)$ fall into the inequality $9b + 108c < 1000 < 11b + 111c$.  

 For $1000 - n = 9k \le 9(7 \cdot 12 + 11) = 855$, we can express $k$ in terms of $(b,c): n \equiv b \pmod{12}, 0 \le b \le 11$ and $c = \frac{n-b}{12} \le 7$ (in other words, we take the greatest possible value of $c$, and then "fill in" the remainder by incrementing $b$). Then $11b + 111c \le 855 + 2b + 3c \le 855 + 2(11) + 3(7) = 898 < 1000$, so these values work. 

 Similarily, for $855 \le 9k \le 9(8 \cdot 12 + 10) = 954$, we can let $(b,c) = (k-8 \cdot 12,8)$, and the inequality $11b + 111c \le 954 + 2b + 3c \le 954 + 2(10) + 3(8) = 998 < 1000$. However, for $9k \ge 963 \Longrightarrow n \le 37$, we can no longer apply this approach. 

 So we now have to examine the numbers on an individual basis. For $9k = 972$, $(b,c) = (0,9)$ works. For $9k = 963, 981, 990, 999 \Longrightarrow n = 37, 19, 10, 1$, we find (using that respectively, $b = 11,9,10,11 + 12p$ for integers $p$) that their is no way to satisfy the inequality $11b + 111c < 1000$. 

 Thus, the answer is $112 - 4 = \boxed{108}$.

 


 A note: Above, we formulated the solution in a forward manner (the last four paragraphs are devoted to showing that all the solutions we found worked except for the four cases pointed out; in a contest setting, we wouldn't need to be nearly as rigorous). A more natural manner of attacking the problem is to think of the process in reverse, namely seeing that $n \equiv 1 \pmod{9}$, and noting that small values of $n$ would not work.  

 Looking at the number $7000$, we obviously see the maximum number of $7's$: a string of $1000 \ 7's$.  Then, we see that the minimum is $28 \ 7's: \ 777*9 + 7 = 7000$.  The next step is to see by what interval the value of $n$ increases.  Since $777$ is $3 \ 7's, \ 77*10 + 7$ is $21 \ 7's$, we can convert a $777$ into $77's$ and $7's$ and add $18$ to the value of $n$.  Since we have $9 \ 777's$ to work with, this gives us $28,46,64,82,100,118,136,154,172,190 ( = 28 + 18n | 1\leq n\leq 9)$ as values for $n$.  Since $77$ can be converted into $7*11$, we can add $9$ to $n$ by converting $77$ into $7's$.  Our $n = 190$, which has $0 \ 777's \ 90 \ 77's \ 10 7's$.  We therefore can add $9$ to $n \ 90$ times by doing this.  All values of $n$ not covered by this can be dealt with with the $n = 46 \ (8 \ 777's \ 10 \ 77's \ 2 \ 7's)$ up to $190$.

 





\textbf{Solution 2}

To simplify, replace all the $7$’s with $1$’s. Because the sum is congruent to $n \pmod 9$ and 
 \[1000 \equiv 1 \pmod 9 \implies n \equiv 1 \pmod 9\]Also, $n \leq 1000$. There are $\big\lfloor \frac{1000}{9} \big\rfloor + 1 = 112$ positive integers that satisfy both conditions i.e. $\{1, 10, 19, 28, 37, 46, . . . , 1000\}.$ 

 For $n = 1, 10, 19$, the greatest sum that is less than or equal to $1000$ is 
 \[6 \cdot 111 + 1 = 667 \implies 112-3 = 109.\]

 Thus $n \geq 28$ and let $S = \{28, 37, 46, . . . , 1000\}$.

 Note that $n=28$ is possible because $9 \cdot 111+1 \cdot 1 = 1000$.

 When $n = 37$, the greatest sum that is at most $1000$ is $8 \cdot 111+6\cdot 11+1 \cdot 1 = 955$. 

 
All other elements of $S$ are possible because if any element $n$ of $S$ between $46$ and $991$ is possible, then $(n+ 9)$ must be too. 

 $\textbf{Case 1:}$ Sum has no $11$'s

 It must have at least one $1$. If it has exactly one $1$, there must be nine $111$’s and $n = 28$. Thus, for $n \geq 46$, the sum has more than one $1$, so it must have at least $1000 - 8 \cdot 111 = 112$ number of $1$’s.For $n \leq 1000$, at least one $111$. To show that if $n$ is possible, then $(n + 9)$ is possible, replace a $111$ with $1 + 1 + 1$, replace eleven $(1 + 1)$’s with eleven $11$’s, and include nine new $1$’s as $+1$’s. The sum remains $1000$.

 $\textbf{Case 2:}$ Sum has at least one $11$.

 Replace an $11$ with $1 + 1$, and include nine new $1$’s as $+1$’s.Now note that $46$ is possible because $8 \cdot 111 + 10 \cdot 11 + 2 \cdot 1 = 1000$.Thus all elements of $S$ except $37$ are possible. 

 Thus there are $\boxed{108}$ possible values for $n$.

 





\textbf{Solution 3}

It's obvious that we cannot have any number $\ge 7777$ because $7777 > 7000$ so the max number that an occur is $777$

 Let's say we have $a$  777's , $b$  77's and $c$  7's

 From here we get our required equation as $777a + 77b + 7c = 7000$

 Now comes the main problem , one might think that if we find number of $(a,b,c)$ then we're done , but in reality we are over-counting our number of $n$'s. This is because $n$ is the total number of 7's and from our equation we'll get $n$ as $3a + 2b + c$ (because there are three 7's , two 7's and one 7) 

 The reason why we're over-counting is because , say $a_1 , b_1 , c_1$ be a solution of our original equation and $a_2 , b_2 , c_2$ be another solution of our original equation , then there can be a possibility that $3a_1 + 2b_1 + c_1 = 3a_2 + 2b_2 + c_2$ where $a_1 \neq a_2 , b_1 \neq b_2 , c_1 \neq c_2$ (example : $2 + 3 + 4 = 1 + 5 + 3$ but $2 \neq 1 , 3 \neq 5 , 4 \neq 3$

 We know that $0 \le a \le 9$ , $0 \le b \le 90$ , $0 \le c \le 1000$The bound on $a$ is easier to handle with , so lets start putting values on $a$ and calculate $b , c , n$ by making cases

 Reduced equation : $111a + 11b + c = 1000$

 Case 1 : $a = 9$

 We get $11b + c = 1 \implies b = 0 , c = 1$ is our only solution thus only $\boxed{1}$ value of $n$

 Case 2 : $a = 8$

 We get $11b + c = 112 \implies b = 0,1,2,\cdots ,10$ and $c = 112 , 101 , 90,\cdots ,2$ From here we get different $n$'s as $24 + 112 , 24 + 103 , 24 + 94, \cdots ,24 + 22$  (remember $n = 3a + 2b + c$ and if you have difficulty in understanding how we got $n = 24 + (\cdots)$ then just put the values of $a,b,c$ i am sure you will get it :) )

 Let's write the sequences of $n$'s in a compact form , $T_p = 24 + 22 + 9(p-1)$ (This will be helpful later on)There are $\boxed{11}$ values of $n$

 Case 3 : $a = 7$

 We get $11b + c = 223 \implies b=0,1,2, \cdots ,20$ and $c = 223,212,201, \cdots ,3$ From here we get different $n$'s as $21 + 223, 21 + 214 , 21 + 205, \cdots ,21 + 43$

 Let's also write the sequences of $n$'s in a compact form , $T_q = 21 + 43 + 9(q-1)$

 Now comes the major part , since we need to find distinct $n$'s so let's subtract the cases where we find common values , from the total number of values.

 To do this we need to make $T_p = T_q \implies p - q = 2$ (after some calculations you'll get $p - q = 2$) . Now we know that $0 \le p \le 10$ and $0 \le q \le 20$ so we get $p$ as $10,9,\cdots,2$ and $q$ as $8,7,\cdots,0$ so there are 9 common solutions out of 21(diff values of $q$) total , so there are $\boxed{21 - 9}$ values of $n$

 Case 4 : $a = 6$

 We get $11b + c = 334 \implies b = 0,1,2, \cdots ,30$ and $c = 334,323,312, \cdots ,4$ From here we get the different $n$'s as $18 + 334,18 + 325,18 + 316, \cdots ,18 + 64$

 Compact form of terms is $T_r = 18 + 64 + 9(r-1)$Now let's repeat the process of eliminating common solutions (one thing to notice is that we've removed common solutions of case 2 from case 3 so if we check case 4 with case 3 we'll remove common solutions from all previous cases and hence we do not have to handle common solutions of case 4 with case 2,1 and in general case X with case X-2,X-3,...2,1 , {basically means checking case X with case X-1 is enough})

 $T_q = T_r \implies 21 + 43 + 9(q-1) = 18 + 64 + 9(r-1) \implies q - r = 2$$\implies q = 20,19,18, \cdots ,2$ and $r = 18,17,16, \cdots ,0$ so there are 19 common solutions out of 31 (diff values of r) so there are $\boxed{31 - 19}$ values of $n$

 Now hopefully you've seen a pattern , if not try another case 5 with $a = 5$ you'll get $\boxed{41 - 29}$ values of $n$

 Like this we get the value of "distinct" $n$ by summing all the sub-values from the cases that we've solved.

 $n = 1 + 11 + (21 - 9) + (31 - 19) + (41 - 29) + \cdots + (91 - 79)$

 $\implies n = (1 + 11 + 21 + 31 + 41 + \cdots + 91) - (0 + 0 + 9 + 19 + 29 + \cdots + 79)$

 $\implies \boxed{n = 108}$

 


 





\textbf{Solution 4 (Construction, not counting)}

First, we note that by the sun of digits and modulo 9 formula, $n \equiv 1000 \equiv 1 \pmod {9}$. Obviously, $1000 = 111+111+111+111+111+111+111+111+111+1$ with $n=28$ is the smallest possible $n$. Now, $111=10 \cdot 11 +1$, $11= 11 \cdot 1$. The first equation applied on a representation with $n$ 1's (i.e. replacing a $111$, if available, with 10 $11$'s and one 1) will increase $n$ by 18, while the second equation applied onto a representation will increase $n$ by 9 only. Thus, there are representations of $n= 46,55,64, \dots , 1000$, since for $46$ we can just replace a $111$ in $n=28$ representation with 10 $11$'s and one 1, and for the rest we use the first equation $111= 10 \cdot 11 +1$ to "increase the range" and the second equation to "fill in the gaps" of the first equation. Essentially we use the first equation for every other multiple of nine and the second for the other multiples of nine and for once the $111$'s run out. Next, obviously $n=1000$ is the maximum $n$ one can achieve. Thus, the solutions $n=28,46,55,64, \dots ,1000$ is all the solutions (one can check there are $\boxed{108}$) if we can show $37$ is impossible. But it is, since $a+b+c=37, 111a+11b+c=1000$ have no Diophantine solution, so we're done.

 
<i>Alternate solutions are always welcome. If you have a different, elegant solution to this problem, please \href{https://artofproblemsolving.com/wiki/index.php?title=2004\_AIME\_II\_Problems/Problem\_14&action=edit}{add it} to this page.</i>

 




	\item (\textit{AIME II 2004}) A long thin strip of paper is 1024 units in length, 1 unit in width, and is divided into 1024 unit squares. The paper is folded in half repeatedly. For the first fold, the right end of the paper is folded over to coincide with and lie on top of the left end. The result is a 512 by 1 strip of double thickness. Next, the right end of this strip is folded over to coincide with and lie on top of the left end, resulting in a 256 by 1 strip of quadruple thickness. This process is repeated 8 more times. After the last fold, the strip has become a stack of 1024 unit squares. How many of these squares lie below the square that was originally the 942nd square counting from the left?


 
 \textbf{Answer: }593



\textbf{Solution 1}

Number the squares $0, 1, 2, 3, ... 2^{k} - 1$.  In this case $k = 10$, but we will consider more generally to find an inductive solution.  Call $s_{n, k}$ the number of squares below the $n$ square after the final fold in a strip of length $2^{k}$.

 Now, consider the strip of length $1024$.  The problem asks for $s_{941, 10}$.  We can derive some useful recurrences for $s_{n, k}$ as follows:  Consider the first fold.  Each square $s$ is now paired with the square $2^{k} - s - 1$.  Now, imagine that we relabel these pairs with the indices $0, 1, 2, 3... 2^{k - 1} - 1$ - then the $s_{n, k}$ value of the pairs correspond with the $s_{n, k - 1}$ values - specifically, double, and maybe $+ 1$ (if the member of the pair that you're looking for is the top one at the final step).

 So, after the first fold on the strip of length $1024$, the $941$ square is on top of the $82$ square.  We can then write

 
 \[s_{941, 10} = 2s_{82, 9} + 1\]

 (We add one because $941$ is the odd member of the pair, and it will be on top.  This is more easily visually demonstrated than proven.)  We can repeat this recurrence, adding one every time we pair an odd to an even (but ignoring the pairing if our current square is the smaller of the two):

 
 \[s_{82, 9} = 2s_{82, 8} = 4s_{82, 7} = 8s_{127 - 82, 6} = 8s_{45, 6}\]

 
 \[s_{45, 6} = 2s_{63 - 45, 5} + 1 = 2s_{18, 5} + 1 = 4s_{31 - 18, 4} + 1 = 4s_{13, 4} + 1\]

 
 \[s_{13, 4} = 2s_{15 - 13, 3} + 1 = 2s_{2, 3}+1\]

 We can easily calculate $s_{2, 3} = 4$ from a diagram.  Plugging back in,

 
 \begin{align*} s_{13, 4} &= 9 \\ s_{45, 6} &= 37 \\ s_{82, 9} &= 296 \\ s_{941, 10} &= \boxed{593}\end{align*}


 





\textbf{Solution 2}

More brute force / thinking about the question logically. If the number doesn't change position, then the number of squares below it does not change. Otherwise, it changes position ($p\mapsto 2^k+1-p$ for the $k$-th fold). We just take the number of squares under it before we folded and now these are above the square. 

 Initially it is in position $942$ with $0$ squares below it. 

 
 \[\begin{array}{llll} \text{\bf Fold 1.} & 942 \text{ is to the right.} & 942 \mapsto 83 & \text{There is } 2^1 - 1 - 0=\boxed{1} \text{ square below it.} \\  &&& \\ \textbf{Fold 2.} & 83 \text{ is to the left.} & 83 \mapsto 83 & \text{There is } \boxed{1} \text{ square below it.} \\ &&& \\ \textbf{Fold 3.} & 83 \text{ is to the left.} & 83 \mapsto 83 & \text{There is } \boxed{1} \text{ square below it.} \\ &&& \\ \text{\bf Fold 4.} & 83 \text{ is to the right.} & 83 \mapsto 46 & \text{There are } 2^4 - 1 - 1 = \boxed{14} \text{ squares below it.}\\  &&& \\ \text{\bf Fold 5.} & 46 \text{ is to the right.} & 46 \mapsto 19 & \text{There are } 2^5 - 1 - 14 = \boxed{17} \text{ squares below it.}\\  &&& \\ \text{\bf Fold 6.} & 19 \text{ is to the right.} & 19 \mapsto 14 & \text{There are } 2^6 - 1 - 17 = \boxed{46} \text{ squares below it.}\\  &&& \\ \text{\bf Fold 7.} & 14 \text{ is to the right.} & 14 \mapsto 3 & \text{There are } 2^7 - 1 - 46 = \boxed{81} \text{ squares below it.}\\  &&& \\ \textbf{Fold 8.} & 3 \text{ is to the left.} & 3 \mapsto 3 & \text{There are } \boxed{81} \text{ squares below it.} \\ &&& \\ \text{\bf Fold 9.} & 3 \text{ is to the right.} & 3 \mapsto 2 & \text{There are } 2^9 - 1 - 81  = \boxed{430} \text{ squares below it.}\\  &&& \\ \text{\bf Fold 10.} & 2 \text{ is to the right.} & 2 \mapsto 1 & \text{There are } 2^{10} - 1 - 430 = \boxed{593} \text{ squares below it.}\\  \end{array}\]

 





\textbf{Solution 3}

We can keep track of the position of the square labeled 942 in each step. We use an $(x,y)$ coordinate system, so originally the 942 square is in the position $(942,1)$. In general, suppose that we've folded the strip into an array $r=2^k$ squares wide and $c=1024/r=2^{10-k}$ squares tall (so we've made $10-k$ folds). Then if a square occupies the location $(x,y)$, we find that after the next fold, it will be in the location described by the procedure
 \[(x,y)\to\begin{cases}(x,y)&\text{if }x\le 2^{k-1}\\ (r+1-x,2c+1-y)&\text{otherwise}.\end{cases}\]Therefore, we can keep track of the square's location in the following table.
 \[\begin{array}{c|c|c} (x,y)&\text{rows}&\text{columns}\\\hline (942,1)&1024&1\\ (83,2)&512&2\\ (83,2)&256&4\\ (83,2)&128&8\\ (46,15)&64&16\\ (19,18)&32&32\\ (14,47)&16&64\\ (3,82)&8&128\\ (3,82)&4&256\\ (2,431)&2&512\\ (1,594)&1&1024.\\ \end{array}\]Therefore, at the end of the process, the square labeled 942 will be in the position $(1,594)$, i.e., it will be above $\boxed{593}$ squares.

 




\end{enumerate}